\documentclass[10pt,a4paper]{amsart}
\usepackage[margin=19mm]{geometry}
\usepackage{amsmath,amssymb,mathtools}
\usepackage[scr=rsfs,cal=euler]{mathalfa}
\usepackage{xfrac}
\usepackage[unicode,hidelinks,bookmarks=false]{hyperref}
\newcommand{\ie}{i.e.\ }
\newcommand{\cf}{cf.\ }
\newcommand{\resp}{respectively\ }
\newcommand{\Subsection}{\subsection}
\providecommand{\nobreakdash}{\nobreak\hbox{-}}
\setlength{\emergencystretch}{2em}
\multlinegap0pt
\usepackage{amssymb}
\usepackage{enumerate}
\usepackage{xcolor}
\usepackage{tikz}
\newtheorem{lemma}{Lemma}
\newtheorem{proposition}{Proposition}
\renewcommand{\H}{\mathbb{H}}
\newcommand{\N}{\mathbb{N}} %este lo añadi yo
\title{On the non-expansiveness of the geodesic flow on surfaces with cusps}
\author{Sergi Burniol Clotet, Fran\c{c}oise Dal'Bo}
\date{Source: arXiv:2603.24310v2 (CC BY 4.0)}
\begin{document}
\maketitle

\noindent\emph{Source and licence.} On the non-expansiveness of the geodesic flow on surfaces with cusps. Version arXiv:2603.24310v2 (CC BY 4.0), distributed by arXiv under the Creative Commons Attribution 4.0 International licence, http://creativecommons.org/licenses/by/4.0/, as recorded in the arXiv licence metadata for this exact version and shown on the abstract page https://arxiv.org/abs/2603.24310v2. Full licence notice: \texttt{licenses/arxiv-2603.24310-CC-BY-4.0.txt}.\par\smallskip

\noindent\textit{Editorial scope.} Counted unit: the article's proposition Pm (source0.tex lines 910--922). The article's complete original proof of it is delivered with it (source0.tex lines 924--989, one \texttt{proof} environment of 66 lines). No mathematics is written, repaired or re-derived: the statement and the proof are the article's own lines, in the article's own notation, and no retained line is substituted or re-flowed.  Retained uncounted as the source-local context the proof needs: lines 309--314; lines 435--445; lines 868--872.  Nothing was omitted from any retained span, so the package carries no omission record.  No retained line contains a citation, so the wrapper carries no bibliography and no bibliography entry.  No retained line contains a figure environment, an image-inclusion command or drawing code, so no image is delivered and the package declares no asset dependency.  Editorial formatting only, in addition: an AMS \texttt{amsart} preamble that restates the article's own declarations and macros used by the retained lines, namely the environments \texttt{lemma}, \texttt{proposition} on separate counters, together with the standard packages those lines need.  The retained environments print in this document as Lemma 1, Proposition 1 in order of appearance, whereas the article prints the same items with the numbering of its own full document.  The complete native source of the article as distributed in the arXiv e-print for this exact version is bundled unchanged as \texttt{sources/197136/source0.tex}.
	\begin{lemma}\label{distance_horo}
		Let $u,v\in T^1\H$ such that $v\in H^{ss} u$. Then for all $t\ge0$
		\[
		\sinh\left(\frac{d(v(t),u(t))}{2}\right) = \sinh\left(\frac{d(v(0),u(0))}{2}\right)e^{-t}.
		\]
	\end{lemma}

 \begin{proposition}[Key Proposition]\label{key_prop} 
 	Let $\tilde u\in T^1 \H$ be tangent to an oriented pair $(\tilde H, p)$.  	
 	Then, for every $t\ge 0$,
 	\begin{align*}
 		\min \{ d_1\!\left(g_{t+\tau_{\tilde H, p ,\tilde u}}\bigl(\operatorname{Wind}_{(\tilde H, p )}(\tilde u)\bigr),\, g_t\tilde u\right),
 		d_1\!\left(g_{t+\tau_{\tilde H, p ,\tilde u}}\bigl(\operatorname{Wind}_{(\tilde H, p )}(\tilde u)\bigr),\, p g_t\tilde u\right)
 		\}
 		\\ \le 12 \ell(\tilde H, p).
 	\end{align*}
 	
 \end{proposition}

\begin{proposition}\label{convergence_winding_time}
	Let $(\alpha_n)_{n\in\\N}$ be a subsequence of $(p_n)_{n\in \N}$ such that \[ \sum_{n=0}^{+\infty} \ell_n^{\alpha} < +\infty.\] Then the sequence $(r_n)_{n\in \mathbb{N}}$ converges towards a real number $r_{\alpha}$. 
	Moreover, $|r_n|\le \sum_{i=0}^{+\infty} \ell_i^{\alpha}$ for all $n\ge 0$ and, hence, $|r_{\alpha}|\le \sum_{i=0}^{+\infty} \ell_i^{\alpha}$ 
	%sum of l_n convergent is enough, generalize
\end{proposition}

\begin{proposition}
	Let $\varepsilon>0$ and let $(\alpha_n)$ be a subsequence of $(p_n)$ such that
	\[
	\ell_n^\alpha < \frac{1}{12}\frac{\varepsilon}{4^n}.
	\]
	Then, for every $l\in \N$, there exists a time $T_l \ge 0$ such that, for all
	$n \in \mathbb{N}$,
	\begin{equation}\label{Pm}
			\forall t\ge T_l,\quad d_1\big(g_{t+r_n} v_n,\, g_t u\big)
		<
		\left(\sum_{k=0}^{n} \frac{1}{2^k}\right)\frac{\varepsilon}{2^l}.
	\end{equation}
\end{proposition}

\begin{proof}
	We first observe that $\sum_{i=0}^{+\infty}\ell^{\alpha}_{i}\le \varepsilon/9$, so we are in the hypothesis of Proposition \ref{convergence_winding_time}.
	
	We start by defining the times $T_l$. 
	Recall that $\beta_n^{-1}\tilde v_n$ is a vector pointing at $\infty$
	and $\beta_n^{-1} g_{r_n}\tilde v_n \in H^{ss} \tilde u$.
	Let
	\[
	D_n = \max_{0 \le k \le n}
	d\big(\beta_k^{-1} g_{r_k}\tilde v_k(0),\, i\big).
	\]
	
	We now set $T_0 = \varepsilon/ 9$ and, for $l \ge 1$,
	\[
	T_l
	= \max \left(
	\log\!\left(
	\sinh\!\left(\frac{D_{l-1}}{2}\right)\cdot 2^{l+2} \cdot \varepsilon^{-1}
	\right), \frac{\varepsilon}{9} \right)
	\]
	
	We will show that this sequence of times satisfies the property \ref{Pm}
	of the statement.
	The proof is by induction. We say that the property $P_m$ is satisfied
	if \ref{Pm} is satisfied for all $n,l \in \{0,\dots,m\}$.
	The statement follows if we prove $P_m$ for all $m$.
	
	We start by showing that $P_0$ is satisfied.
	Recall that $\tilde v_0 = \mathrm{Wind}_{(\tilde H^{\alpha}_0, \alpha_0)}(\tilde u)$. By Proposition \ref{key_prop}, then we have, for all $t\ge 0$
	\[
	d_1(g_{t+r_0}v_0, g_t u) \le  12 \ell^{\alpha}_0 < \varepsilon.
	\]
	
	
	We now assume that $P_m$, $m\ge0$, is satisfied and show $P_{m+1}$. We need to show \ref{Pm} for $n,l\in \{0,\dots,m+1\}$. If we take $n,l\in \{0,\dots,m\}$ then \ref{Pm} is satisfied by the induction assumption.
	
	Let us first consider the case $n \in \{0,\dots,m\}$ and $l=m+1$. Observe that $g_{r_n} \beta_n^{-1}\tilde v_n$ belongs to the horocycle $H^{ss}\tilde u$. Then we can estimate the distance $d_1$ by the distance between the basepoints and use Lemma \ref{distance_horo}, after applying the inequality $x\le \sinh(x)$ for $x\ge0$. From the definition of the time $T_{m+1}$ we obtain that, for all $t\ge T_{m+1}$, 
	\begin{align*}
		d_1\big(g_{t+r_n} \beta_n^{-1} \tilde v_n,\, g_t \tilde u\big) \le & 
		2 d\big( \beta_n^{-1} \tilde v_n (t+r_n),\, \tilde u (t) \big) \le 
		4 \sinh (\frac{d( \beta_n^{-1} \tilde v_n (t+r_n),\, \tilde u (t) )}{2}) \\
		\le & 4 \sinh (\frac{d( \beta_n^{-1} \tilde v_n (r_n),\, \tilde u (0) )}{2}) e^{-t} \le 4 \sinh(\frac{D_m}{2}) e^{-t}	\le \frac{\varepsilon}{2^{m+1}}.	
	\end{align*}
Hence, for all $t\ge T_{m+1}$,
\[
d_1\big(g_{t+r_n} v_n,\, g_t u\big) \le \frac{\varepsilon}{2^{m+1}} \le \left(\sum_{k=0}^{n} \frac{1}{2^k}\right)\frac{\varepsilon}{2^{m+1}}.
\]
So far we have proved that \ref{Pm} is verified for any $n\in\{0,\dots,m\}$ and any $l\in \{0,\dots,m+1\}$. 

Finally, let us consider the case that $n=m+1$ and $l\in \{0,\dots,m+1\}$. 
Recall that \[\tilde v_{m+1}=\beta_m \mathrm{Wind}_{(\tilde  H'_{m+1}, \alpha_{m+1})}(\beta_m^{-1} \tilde v_m)\] and $\tau_{\tilde H_{m+1}', \alpha_{m+1}, \beta_{m}^{-1}\tilde v_m} = r_{m+1}-r_m$. Applying Proposition \ref{key_prop}, we have
\[
d_1\big(g_{s+r_{m+1}-r_m} \beta_m^{-1}\tilde v_{m+1},\, g_{s} \beta_m^{-1} \tilde v_m\big) \le 12 \ell^{\alpha}_{m+1} < \frac{\varepsilon}{4^{m+1}}
\]
for all $s\ge 0$. Hence,
\[
d_1\big(g_{t+r_{m+1}} v_{m+1},\, g_{t+r_m} v_m\big) < \frac{\varepsilon}{4^{m+1}}
\]
for all $t\ge r_m$, and in particular this holds if $t\ge \varepsilon/ 9$ thanks to Proposition \ref{convergence_winding_time}. Using \ref{Pm}, that we already established for $n\in\{0,\dots,m\}$ and any $l\in \{0,\dots,m+1\}$, if $t\ge T_l$, we have
\begin{align*}
	d_1\big(g_{t+r_{m+1}} v_{m+1},\, g_t u\big) \le & \, d_1\big(g_{t+r_{m+1}} v_{m+1},\, g_{t+r_m} v_m\big) + d_1\big(g_{t+r_m} v_m,\, g_t u\big) \\
	 <& \, \frac{\varepsilon}{4^{m+1}} + \left(\sum_{k=0}^{m} \frac{1}{2^k}\right)\frac{\varepsilon}{2^l} \le \left(\sum_{k=0}^{m+1} \frac{1}{2^k}\right)\frac{\varepsilon}{2^l}.
\end{align*}

	
\end{proof}

\end{document}
